\begin{figure}[H]
	\centering
	\includegraphics[width=1\textwidth]{SpecificaTecnica/Chat/VecDbBackend}
	\caption{Diagramma delle classi per il modulo di ricerca vettoriale: Backend}
\end{figure}

Per la definizione di ToolCatalogService, vedere la sezione \ref{ToolCatalogService}.\\
Per la definizione di \textit{EmbeddedCartService}, vedere \ref{EmbeddedCartService}.\\

\subsubsubsection{FaissCartDb}

\begin{figure}[H]
	\centering
	\includegraphics[width=0.7\textwidth]{SpecificaTecnica/Chat/FaissCartDb}
	\caption{Ricerca Vettoriale: FaissCartDb}
\end{figure}

\paragraph{Descrizione}
Classe che gestisce un indice FAISS$^G$ per la ricerca vettoriale nel carrello.\\
Mantiene una mappa \textit{user\_map} che associa ogni username alle posizioni dei propri prodotti nell'indice, in modo da filtrare i risultati in fase di ricerca.

\paragraph{Attributi}
\begin{itemize}
	\item \textbf{dimension: int} - dimensione dei vettori dell'indice.
	\item \textbf{index: faiss.IndexFlatL2} - indice FAISS$^G$ basato sulla distanza euclidea L2.
	\item \textbf{prod\_ids: List[string]} - lista degli identificativi di prodotto ordinata per posizione nell'indice.
	\item \textbf{user\_map: Dict[string, List[int]]} - dizionario che mappa ogni username alla lista delle posizioni nell'indice occupate dai suoi prodotti.
\end{itemize}

\paragraph{Metodi e funzioni}
\begin{itemize}
	\item \textbf{\_\_init\_\_(dimension: int)} - inizializza il database$^G$ FAISS$^G$ con la dimensione specificata, creando l'indice e le strutture dati necessarie.
	\item \textbf{add(prod\_id: string, username: string, vector: np.ndarray)} - aggiunge al database$^G$ vettoriale il prodotto identificato da \textit{prod\_id} per l'utente \textit{username} con il relativo vettore.
	\item \textbf{search(vector: np.ndarray, n: int, threshold: float, username: string): List[string]} - ricerca nell'indice i vettori più vicini a \textit{vector}, filtrando i risultati per le sole posizioni associate a \textit{username} e restituendo gli identificativi con distanza L2 inferiore o uguale a \textit{threshold}.
\end{itemize}

\subsubsubsection{VecDbPortOut}

\begin{figure}[H]
	\centering
	\includegraphics[width=0.6\textwidth]{SpecificaTecnica/Chat/VecDbPortOut}
	\caption{Ricerca Vettoriale: VecDbPortOut}
\end{figure}

\paragraph{Descrizione}
Classe astratta che definisce la porta di output verso il database$^G$ vettoriale.\\
Disaccoppia i servizi applicativi (\textit{EmbeddedCartService},
\textit{EmbeddedCatalogService}) dalla specifica implementazione del database$^G$ vettoriale.\\
Le classi adapter concrete (\textit{CartVecDbAdapter}, \textit{CatalogVecDbAdapter})
devono implementare tutti i metodi qui dichiarati.

\paragraph{Attributi}
Non definisce attributi propri, in quanto funge esclusivamente da interfaccia.

\paragraph{Metodi e funzioni}
\begin{itemize}
	\item \textbf{reset()} - svuota completamente il database$^G$ vettoriale, eliminando tutti
	i vettori e le associazioni precedentemente inserite.
	\item \textbf{search(v: np.ndarray, n: int, threshold: float): List[string]} -
	ricerca nel database$^G$ vettoriale i prodotti il cui vettore è simile a \textit{v}
	entro la soglia indicata, restituendo al massimo \textit{n} identificativi.
	\item \textbf{add(prod\_id: string, vector: np.ndarray)} - aggiunge al database$^G$
	vettoriale il prodotto identificato da \textit{prod\_id} con il relativo vettore.
\end{itemize}

\subsubsubsection{CartVecDbAdapter}

\begin{figure}[H]
	\centering
	\includegraphics[width=0.6\textwidth]{SpecificaTecnica/Chat/CartVecDbAdapter}
	\caption{Ricerca Vettoriale: CartVecDbAdapter}
\end{figure}

\paragraph{Descrizione}
Classe che implementa \textit{VecDbPortOut} seguendo il pattern degli adattatori,
specializzandosi per la gestione del database$^G$ vettoriale del carrello.\\
Riceve alla costruzione un'istanza di \textit{FaissCartDb} e il nome utente corrente, e
trasparisce le operazioni sul database$^G$ delegando a \textit{FaissCartDb} includendo il
contesto utente.

\paragraph{Attributi}
\begin{itemize}
	\item \textbf{faiss\_db: FaissCartDb} - istanza del database$^G$ vettoriale FAISS$^G$
		dedicato al carrello.
	\item \textbf{username: string} - nome utente associato alla sessione corrente,
		utilizzato per filtrare i risultati della ricerca.
\end{itemize}

\paragraph{Metodi e funzioni}
Oltre alle implementazioni concrete di \textit{VecDbPortOut}, è disponibile la seguente funzione:
\begin{itemize}
	\item \textbf{\_\_init\_\_(faiss\_db: FaissCartDb, username: string)} -
		inizializza l'adattatore ricevendo il database$^G$ FAISS$^G$ e il nome utente.
\end{itemize}

\subsubsubsection{FaissCatalogDb}

\begin{figure}[H]
	\centering
	\includegraphics[width=0.7\textwidth]{SpecificaTecnica/Chat/FaissCatalogDb}
	\caption{Ricerca Vettoriale: FaissCatalogDb}
\end{figure}

\paragraph{Descrizione}
Classe che gestisce un indice FAISS$^G$ per la ricerca vettoriale sul catalogo
prodotti.\\
A differenza di \textit{FaissCartDb} non mantiene una mappa per utente, in quanto il
catalogo è condiviso tra tutti gli utenti del sistema.

\paragraph{Attributi}
\begin{itemize}
	\item \textbf{dimension: int} - dimensione dei vettori dell'indice, impostata a 284
		per default.
	\item \textbf{index: faiss.IndexFlatL2} - indice FAISS$^G$ basato sulla distanza
		euclidea L2.
	\item \textbf{prod\_ids: List[string]} - lista degli identificativi di prodotto
		ordinata per posizione nell'indice.
\end{itemize}

\paragraph{Metodi e funzioni}
\begin{itemize}
	\item \textbf{\_\_init\_\_(dimension: int)} - inizializza il database$^G$ FAISS$^G$ con la dimensione specificata, creando l'indice e le strutture dati necessarie.
	\item \textbf{search(vector: np.ndarray, n: int, threshold: float): List[string]} -
		ricerca nell'indice i vettori più vicini a \textit{vector} e restituisce gli
		identificativi con distanza L2 inferiore o uguale a \textit{threshold}.
\end{itemize}

\subsubsubsection{CatalogVecDbAdapter}

\begin{figure}[H]
	\centering
	\includegraphics[width=0.6\textwidth]{SpecificaTecnica/Chat/CatalogVecDbAdapter}
	\caption{Ricerca Vettoriale: CatalogVecDbAdapter}
\end{figure}

\paragraph{Descrizione}
Classe che implementa \textit{VecDbPortOut} seguendo il pattern degli adattatori,
specializzandosi per la gestione del database$^G$ vettoriale del catalogo.\\
Riceve alla costruzione un'istanza di \textit{FaissCatalogDb} e delega ad essa tutte le
operazioni senza aggiungere il contesto utente, in quanto il catalogo è condiviso.

\paragraph{Attributi}
\begin{itemize}
	\item \textbf{faiss\_db: FaissCatalogDb} - istanza del database$^G$ vettoriale FAISS$^G$
		dedicato al catalogo.
\end{itemize}

\paragraph{Metodi e funzioni}
Oltre alle implementazioni concrete di \textit{VecDbPortOut}, è disponibile la seguente funzione:
\begin{itemize}
	\item \textbf{\_\_init\_\_(faiss\_db: FaissCatalogDb)} - inizializza l'adattatore
		ricevendo il database$^G$ FAISS$^G$ del catalogo.
\end{itemize}

\subsubsubsection{EmbedderPort}

\begin{figure}[H]
	\centering
	
\includegraphics[width=0.5\textwidth]{SpecificaTecnica/Chat/EmbedderPort}
	\caption{Ricerca Vettoriale: EmbedderPort}
\end{figure}

\paragraph{Descrizione}
Classe astratta che definisce la porta per la generazione di embedding testuali.\\
Segue il pattern architetturale delle porte e degli adattatori, disaccoppiando la logica
applicativa di \textit{EmbeddedCartService} ed \textit{EmbeddedCatalogService} dalla
specifica libreria di embedding utilizzata.\\
Tutte le classi adapter concrete devono implementare il metodo qui dichiarato.

\paragraph{Attributi}
Non definisce attributi propri, in quanto funge esclusivamente da interfaccia.

\paragraph{Metodi e funzioni}
\begin{itemize}
	\item \textbf{embed(text: string): np.ndarray} - restituisce il vettore di embedding
	corrispondente al testo passato in ingresso.
\end{itemize}

\subsubsubsection{EmbedderAdapter}

\begin{figure}[H]
	\centering
	\includegraphics[width=0.8\textwidth]{SpecificaTecnica/Chat/EmbedderAdapter}
	\caption{Ricerca Vettoriale: EmbedderAdapter}
\end{figure}

\paragraph{Descrizione}
Classe che implementa \textit{EmbedderPort} seguendo il pattern degli adattatori.\\
Riceve un'istanza di \textit{SentenceTransformerEmbedder} alla costruzione e delega
interamente ad essa la generazione degli embedding.\\
Permette di sostituire la libreria di embedding sottostante senza modificare i servizi
applicativi.

\paragraph{Attributi}
\begin{itemize}
	\item \textbf{sentence\_transformer\_embedder: SentenceTransformerEmbedder} -
		istanza del componente concreto che esegue la generazione dell'embedding.
\end{itemize}

\paragraph{Metodi e funzioni}
Oltre alle implementazioni concrete di \textit{EmbedderPort}, è disponibile la seguente funzione:
\begin{itemize}
	\item \textbf{\_\_init\_\_(sentence\_transformer\_embedder: SentenceTransformerEmbedder)}
		- inizializza l'adattatore ricevendo l'istanza di
		\textit{SentenceTransformerEmbedder}.
\end{itemize}

\subsubsubsection{SentenceTransformerEmbedder}

\begin{figure}[H]
	\centering
	\includegraphics[width=0.5\textwidth]{SpecificaTecnica/Chat/SentenceTransformerEmbedder}
	\caption{Ricerca Vettoriale: SentenceTransformerEmbedder}
\end{figure}

\paragraph{Descrizione}
Classe che incapsula il modello \textit{all-MiniLM-L6-v2} di
\textit{sentence-transformers} per la generazione di embedding testuali normalizzati.\\
Viene istanziata una sola volta e condivisa tramite \textit{EmbedderAdapter}.\\
Non implementa direttamente \textit{EmbedderPort}: tale responsabilità è delegata ad
\textit{EmbedderAdapter}.

\paragraph{Attributi}
\begin{itemize}
	\item \textbf{embedder: SentenceTransformer} - istanza del modello
	\textit{all-MiniLM-L6-v2} caricata al momento della costruzione.
\end{itemize}

\paragraph{Metodi e funzioni}
\begin{itemize}
	\item \textbf{embed(text: string): NDArray[np.float32]} - genera l'embedding
	normalizzato del testo, convertito in un array NumPy di tipo \textit{float32}.
\end{itemize}

\subsubsubsection{EmbeddedCatalogService}

\begin{figure}[H]
	\centering
	\includegraphics[width=0.8\textwidth]{SpecificaTecnica/Chat/EmbeddedCatalogService}
	\caption{Ricerca Vettoriale: EmbeddedCatalogService}
\end{figure}
\label{EmbeddedCatalogService}

\paragraph{Descrizione}
Classe che coordina la ricerca vettoriale sull'intero catalogo prodotti.\\
Espone i metodi per caricare il catalogo nel database$^G$ vettoriale e per cercare prodotti
semanticamente simili a una query$^G$ testuale.

\paragraph{Attributi}
\begin{itemize}
	\item \textbf{catalog\_vect: VecDbPortOut} - istanza del porto di output verso il
		database$^G$ vettoriale del catalogo.
	\item \textbf{catalog\_repo: CatalogRepoPort} - istanza del porto di accesso ai
		dati del catalogo.
	\item \textbf{embedder: EmbedderPort} - istanza del porto per la generazione degli
		embedding.
\end{itemize}

\paragraph{Metodi e funzioni}
\begin{itemize}
	\item \textbf{\_\_init\_\_(catalog\_vect: VecDbPortOut, catalog\_repo: CatalogRepoPort, embedder: EmbedderPort)} - inizializza il servizio ricevendo le dipendenze tramite iniezione.
	\item \textbf{search\_catalog(query$^G$: string, threshold: float): List[string]} - genera l'embedding della query$^G$ e restituisce gli identificativi dei prodotti del catalogo semanticamente simili entro la soglia indicata.
	\item \textbf{load\_catalog()} - recupera l'intero catalogo tramite \textit{catalog\_repo}, genera l'embedding di ciascun prodotto e lo aggiunge al database$^G$ vettoriale.
\end{itemize}

%-------------------------------------------------------------------------------------------------------------
\begin{figure}[H]
	\centering
	\includegraphics[width=0.7\textwidth]{SpecificaTecnica/Chat/CatalogBackend}
	\caption{Diagramma delle classi per il catalogo: Backend}
\end{figure}
Per la definizione di ToolCatalogService, vedere \ref{ToolCatalogService}

\subsubsubsection{CatalogRepoPort}

\begin{figure}[H]
	\centering
	\includegraphics[width=0.6\textwidth]{SpecificaTecnica/Chat/CatalogRepoPort}
	\caption{Catalogo: CatalogRepoPort}
\end{figure}
\label{CatalogRepoPort}

\paragraph{Descrizione}
Classe astratta che definisce la porta di accesso ai dati del catalogo prodotti.\\
Disaccoppia i servizi applicativi (\textit{EmbeddedCatalogService},
\textit{ToolCatalogService}) dalla specifica implementazione del repository.\\
Tutte le classi adapter concrete devono implementare i metodi qui dichiarati.

\paragraph{Attributi}
Non definisce attributi propri, in quanto funge esclusivamente da interfaccia.

\paragraph{Metodi e funzioni}
\begin{itemize}
	\item \textbf{get\_product(prod\_id: string): CatalogProduct} -
		restituisce il prodotto corrispondente all'identificativo indicato, oppure
		\textit{None} se non trovato.
	\item \textbf{get\_full\_catalog(): List[CatalogProduct]} - restituisce la lista
		completa dei prodotti presenti nel catalogo.
\end{itemize}

\subsubsubsection{CatalogRepoAdapter}

\begin{figure}[H]
	\centering
	\includegraphics[width=0.6\textwidth]{SpecificaTecnica/Chat/CatalogRepoAdapter}
	\caption{Catalogo: CatalogRepoAdapter}
\end{figure}

\paragraph{Descrizione}
Classe che implementa \textit{CatalogRepoPort} seguendo il pattern degli
adattatori.\\
Riceve alla costruzione un'istanza di \textit{CatalogRepository} e si occupa di
convertire i modelli ORM \textit{Anaart} negli schemi applicativi \textit{CatalogProduct},
mappando in particolare l'unità di misura tramite \textit{MeasureUnitEnum}.

\paragraph{Attributi}
\begin{itemize}
	\item \textbf{repo: CatalogRepository} - istanza del repository$^G$ che accede
		direttamente al database.
\end{itemize}

\paragraph{Metodi e funzioni}
Oltre alle implementazioni concrete di \textit{CatalogRepoPort}, sono disponibili le seguenti funzioni:
\begin{itemize}
	\item \textbf{\_\_init\_\_(repo: CatalogRepository)} - inizializza l'adattatore
		ricevendo il repository.
	\item \textbf{\_map\_measure\_unit(row: Anaart): MeasureUnitEnum} - metodo privato
		che converte il codice e la descrizione dell'unità di misura presenti nel record
		ORM nel corrispondente valore dell'enumerazione \textit{MeasureUnitEnum}.
	\item \textbf{\_map\_product(row: Anaart): CatalogProduct} - metodo privato che
		converte un record ORM \textit{Anaart} nello schema applicativo
		\textit{CatalogProduct}, delegando la mappatura dell'unità di misura a
		\textit{\_map\_measure\_unit}.
\end{itemize}

\subsubsubsection{CatalogRepository}

\begin{figure}[H]
	\centering
	\includegraphics[width=0.6\textwidth]{SpecificaTecnica/Chat/CatalogRepository}
	\caption{Catalogo: CatalogRepository}
\end{figure}

\paragraph{Descrizione}
Classe che gestisce l'accesso diretto al database$^G$ per il catalogo prodotti.\\
Riceve una sessione SQLModel alla costruzione e la utilizza per eseguire le query$^G$ tramite il modello ORM \textit{Anaart}.\\
Rappresenta il livello più basso della catena \textit{EmbeddedCatalogService} → \textit{CatalogRepoAdapter} → \textit{CatalogRepository}.

\paragraph{Attributi}
\begin{itemize}
	\item \textbf{db$^G$: Session} - sessione SQLModel attiva utilizzata per l'esecuzione
		delle query.
\end{itemize}

\paragraph{Metodi e funzioni}
\begin{itemize}
	\item \textbf{\_\_init\_\_(db$^G$: Session)} - inizializza il repository$^G$ ricevendo la
		sessione del database.
	\item \textbf{get\_product(prod\_id: string): Anaart} - recupera il prodotto
		corrispondente all'identificativo indicato, restituendo \textit{None} se non
		trovato.
	\item \textbf{get\_full\_catalog(): List[Anaart]} - recupera tutti i prodotti
		presenti nella tabella \textit{Anaart}.
\end{itemize}

\begin{figure}[H]
	\centering 
	\includegraphics[width=1.0\textwidth]{SpecificaTecnica/Chat/ChatBackend}
	\caption{Diagramma delle classi per la gestione del carrello}
\end{figure}
Per la definizione di \textit{ToolCartService}, vedere \ref{ToolCartService}

\paragraph{Descrizione generale}
Questa parte si occupa della gestione del carrello associato a uno specifico utente.\\
Le classi descritte consentono di recuperare i prodotti presenti nel carrello, aggiungere nuovi articoli, rimuoverli, modificarne la quantit\`a e inviare l'ordine finale.\\
L'architettura segue il pattern Ports and Adapters, separando la logica applicativa dai dettagli di persistenza e integrando un servizio dedicato alla ricerca semantica dei prodotti presenti nel carrello.\\

\subsubsubsection{EmbeddedCartService}
\label{EmbeddedCartService}
\begin{figure}[H]
	\centering 
	\includegraphics[width=0.9\textwidth]{SpecificaTecnica/Chat/EmbeddedCartService}
	\caption{Gestione del carrello: EmbeddedCartService}
\end{figure}

\paragraph{Descrizione}
\textit{EmbeddedCartService} \`e il service applicativo responsabile$^G$ del caricamento del contenuto del carrello in forma vettoriale e della successiva ricerca semantica.\\
Recupera i prodotti del carrello tramite la porta \textit{CartRepoPort}, ne costruisce una rappresentazione vettoriale tramite \textit{EmbedderPort} e interroga il database$^G$ vettoriale attraverso \textit{VecDbPortOut}.

\paragraph{Attributi}
\begin{itemize}
	\item \textbf{cart\_vect: VecDbPortOut} - attributo privato di tipo \textit{VecDbPortOut}.\\
	Porta utilizzata per memorizzare e interrogare la rappresentazione vettoriale dei prodotti presenti nel carrello.
	\item \textbf{cart\_repo: CartRepoPort} - attributo privato di tipo \textit{CartRepoPort}.\\
	Porta utilizzata per recuperare i prodotti presenti nel carrello dell'utente.
	\item \textbf{embedder: EmbedderPort} - attributo privato di tipo \textit{EmbedderPort}.\\
	Porta utilizzata per generare gli embedding dei prodotti da indicizzare nel database$^G$ vettoriale.
\end{itemize}

\paragraph{Metodi e funzioni}
\begin{itemize}
	\item \textbf{\_\_init\_\_(cart\_vect: VecDbPortOut, cart\_repo: CartRepoPort, embedder: EmbedderPort)} - metodo pubblico di inizializzazione.\\
	Riceve e memorizza le dipendenze necessarie per il recupero dei dati del carrello, la generazione degli embedding e l'interrogazione vettoriale.
	\item \textbf{load\_cart(username: string)} - metodo pubblico che recupera dal repository$^G$ i prodotti presenti nel carrello dell'utente identificato da \textit{username} e li carica nel sistema di ricerca vettoriale.\\
	Questo passaggio prepara i dati necessari per successive ricerche semantiche.
	\item \textbf{search\_cart(username: string, query$^G$: string, threshold: float): List[string]} - metodo pubblico che esegue una ricerca semantica all'interno del carrello dell'utente, utilizzando la query$^G$ fornita e una soglia di similarità.\\
	Restituisce gli identificativi o i riferimenti ai prodotti che soddisfano i criteri di ricerca.
\end{itemize}

\subsubsubsection{CartService}
\begin{figure}[H]
	\centering 
	\includegraphics[width=0.7\textwidth]{SpecificaTecnica/Chat/CartService}
	\caption{Gestione del carrello: CartService}
\end{figure}

\paragraph{Descrizione}
\textit{CartService} \`e il service applicativo dedicato alle operazioni tradizionali di gestione del carrello e all'invio dell'ordine.\\
Rappresenta un punto di accesso semplificato alle funzionalit\`a di lettura e modifica$^G$ del carrello, delegando le operazioni alla porta \textit{CartRepoPort}.

\paragraph{Attributi}
\begin{itemize}
	\item \textbf{adapter: CartRepoPort} - attributo privato di tipo \textit{CartRepoPort}.\\
	Porta utilizzata per accedere alle operazioni di persistenza del carrello.
\end{itemize}

\paragraph{Metodi e funzioni}
\begin{itemize}
	\item \textbf{\_\_init\_\_(adapter: CartRepoPort)} - metodo pubblico di inizializzazione.\\
	Riceve e memorizza la porta di accesso al repository$^G$ del carrello.
	\item \textbf{get\_cart\_products(username: string): List[CartProduct]} - metodo pubblico che recupera tutti i prodotti presenti nel carrello dell'utente identificato da \textit{username}.\\
	Restituisce una collezione di oggetti di tipo \textit{CartProduct}.
	\item \textbf{send\_order(username: string)} - metodo pubblico che avvia l'invio dell'ordine relativo al carrello dell'utente specificato.\\
	Delega la persistenza finale dell'operazione alla porta \textit{CartRepoPort}.
\end{itemize}

\subsubsubsection{CartRepoPort}
\begin{figure}[H]
	\centering 
	\includegraphics[width=1.0\textwidth]{SpecificaTecnica/Chat/CartRepoPort}
	\caption{Gestione del carrello: CartRepoPort}
\end{figure}
\label{CartRepoPort}

\paragraph{Descrizione}
\textit{CartRepoPort} \`e un'interfaccia che definisce il contratto per l'accesso e la modifica$^G$ dei dati relativi al carrello.\\
Viene implementata da \textit{CartRepoAdapter} e utilizzata dai service applicativi per mantenere il disaccoppiamento tra logica di business e persistenza.

\paragraph{Attributi}
\textit{CartRepoPort} non definisce attributi.

\paragraph{Metodi e funzioni}
\begin{itemize}
	\item \textbf{get\_products(username: string): List[CartProduct]} - metodo pubblico che recupera i prodotti presenti nel carrello dell'utente identificato da \textit{username}.\\
	Restituisce una collezione di oggetti di tipo \textit{CartProduct}.
	\item \textbf{add\_product(prod\_id: string, username: string, qty: int): CartProduct} - metodo pubblico che aggiunge al carrello dell'utente il prodotto identificato da \textit{prod\_id} nella quantit\`a indicata.\\
	Restituisce il prodotto aggiunto o aggiornato.
	\item \textbf{remove\_product(prod\_id: string, username: string): CartProduct} - metodo pubblico che rimuove dal carrello dell'utente il prodotto identificato da \textit{prod\_id}.\\
	Restituisce il prodotto coinvolto nell'operazione.
	\item \textbf{update\_quantity(prod\_id: string, username: string, qty: int, operation: CartUpdateOperation): CartProduct} - metodo pubblico che aggiorna la quantit\`a di un prodotto gi\`a presente nel carrello.\\
	L'aggiornamento avviene secondo l'operazione specificata da \textit{CartUpdateOperation}.
	\item \textbf{send\_order(username: string): void} - metodo pubblico che conferma e inoltra l'ordine associato al carrello dell'utente.
\end{itemize}

\subsubsubsection{CartRepoAdapter}
\begin{figure}[H]
	\centering 
	\includegraphics[width=0.9\textwidth]{SpecificaTecnica/Chat/CartRepoAdapter}
	\caption{Gestione del carrello: CartRepoAdapter}
\end{figure}

\paragraph{Descrizione}
\textit{CartRepoAdapter} \`e la classe concreta che implementa l'interfaccia \textit{CartRepoPort}.\\
Agisce da adattatore tra i service applicativi e il repository$^G$ concreto \textit{CartRepository}, traducendo le operazioni richieste dal dominio in chiamate al layer di persistenza.

\paragraph{Attributi}
\begin{itemize}
	\item \textbf{repo: CartRepository} - attributo privato di tipo \textit{CartRepository}.\\
	Rappresenta il repository$^G$ concreto utilizzato per eseguire le operazioni sul database.
\end{itemize}

\paragraph{Metodi e funzioni}
Oltre ad essere implementate tutte le funzioni di \textit{CartRepoPort}, \`e disponibile anche la seguente funzione:
\begin{itemize}
	\item \textbf{\_\_init\_\_(repo: CartRepository)} - metodo pubblico di inizializzazione.\\
	Riceve e memorizza il repository$^G$ concreto utilizzato per l'accesso ai dati del carrello.
\end{itemize}

\subsubsubsection{CartRepository}
\begin{figure}[H]
	\centering 
	\includegraphics[width=1.0\textwidth]{SpecificaTecnica/Chat/CartRepository}
	\caption{Gestione del carrello: CartRepository}
\end{figure}

\paragraph{Descrizione}
\textit{CartRepository} \`e la classe che gestisce direttamente le operazioni di lettura e scrittura dei dati del carrello sul database.\\
Utilizza una sessione di database$^G$ per eseguire il recupero dei prodotti, l'aggiunta di nuovi articoli, la rimozione di elementi, l'aggiornamento delle quantit\`a e l'invio finale dell'ordine.

\paragraph{Attributi}
\begin{itemize}
	\item \textbf{db$^G$: Session} - attributo privato di tipo \textit{Session}.\\
	Rappresenta la sessione di database$^G$ utilizzata per eseguire le operazioni di persistenza.
\end{itemize}

\paragraph{Metodi e funzioni}
\begin{itemize}
	\item \textbf{\_\_init\_\_(db$^G$: Session)} - metodo pubblico di inizializzazione.\\
	Riceve una sessione di database$^G$ e la memorizza per l'esecuzione delle operazioni sul carrello.
	\item \textbf{get\_products(username: string): List[CartProductRepository]} - metodo pubblico che recupera dal database$^G$ i prodotti presenti nel carrello dell'utente indicato.\\
	Restituisce una collezione di entit\`a repository$^G$ associate ai prodotti del carrello.
	\item \textbf{add\_product(prod\_id: string, username: string, qty: int): CartProductRepository} - metodo pubblico che inserisce nel database$^G$ un nuovo prodotto nel carrello dell'utente oppure ne aggiorna la presenza secondo la logica prevista.\\
	Restituisce l'entit\`a repository$^G$ risultante.
	\item \textbf{remove\_product(username: string, prod\_id: string): CartProductRepository} - metodo pubblico che rimuove dal database$^G$ il prodotto indicato dal carrello dell'utente.\\
	Restituisce l'entit\`a repository$^G$ coinvolta nell'operazione.
	\item \textbf{update\_quantity(username: string, prod\_id: string, qty: int, operation: CartUpdateOperation): CartProductRepository} - metodo pubblico che aggiorna nel database$^G$ la quantit\`a del prodotto indicato nel carrello dell'utente.\\
	La modifica$^G$ viene applicata in accordo all'operazione specificata.
	\item \textbf{send\_order(username: string)} - metodo pubblico che conclude il processo$^G$ di acquisto inviando l'ordine associato al carrello dell'utente.
\end{itemize}